\section{Details of the precursor study}\label{sec:precursor-details}

This appendix records the design and the archived results of the external
presolve study summarized in Section~\ref{sec:experiments-precursor}. The study
is separate from the prospective comparison; its results are not pooled with
those of Section~\ref{sec:experiments}.

\subsection{Design}

\emph{Reference solutions.} For 275 of the 339 instances the reference solution
was a MINLPLib solution file whose objective matched the library's primal
bound; for 63 it was the Gurobi solution with default settings, which had to be
optimal, match the known value to $10^{-4}$, and be feasible to $10^{-5}$. One
instance had no reference solution.

\emph{Relaxation and OBBT.} The OBBT relaxation was an LP built with Gurobi,
with integrality dropped: the four McCormick inequalities for each distinct
bilinear term, five tangents at equally spaced points and the secant for each
square (with $x_i^2$ replaced by $x_i$ for binary $x_i$), every original row with
its quadratic terms replaced by the lifted variables, and the cutoff row. After
a bound change the affected rows were rewritten in place and the LP was warm
started; presolve was off. An LP reported infeasible produced no tightening; this
occurred for two LPs. Each new bound was relaxed by $10^{-6}(1+|\beta|)$ for a
bound $\beta$ and rounded for integer variables, and it was accepted only if it
reduced the width by at least a relative $10^{-4}$. A round minimized and
maximized every variable occurring in a nonlinear term. The main trajectories
used sequential (Gauss--Seidel) updates and skipped a direction for the rest of a
round when an earlier LP solution of that round already attained its bound.
With sequential updates this filtering is a heuristic: a point from an earlier
LP of the round may violate rows rewritten after later updates, so a skipped
direction could still have tightened. It never produces an invalid bound.
After each round the study recorded $\rho$, the total normalized width after the
round divided by the total before, over the variables whose width shrank by at
least $0.1\%$, where widths are normalized by the box obtained from feasibility-based
tightening.

\emph{Stopping rules.} \texttt{r1} and \texttt{r5} stop after one and five
rounds; \texttt{fp} stops after a round in which no bound changed, after 50
rounds, or after 400 seconds. Because of the margins, the filtering, and the
improvement threshold, a round without change is a numerical stopping event,
not an exact fixed-point certificate. The adaptive rules \texttt{ad0.5} and
\texttt{ad0.8} stop after a round with $\rho>\theta$ for $\theta=0.5$ or $0.8$, after a round without change, or
after 120 seconds of OBBT, and restrict rounds after the first to the variables
that moved in the previous round and their bilinear partners. No time limit
below 400 seconds applied to the first round, even for the adaptive rules, whose
120-second cap was checked only between rounds.

\emph{Pipelines and baselines.} A root-node run of the target solver (seed $0$,
node limit one, 30 seconds) supplied the incumbent and the cutoff
$U=\hat f+10^{-6}\max\{1,|\hat f|\}$. If this run proved optimality (Gurobi
101 instances, SCIP 66), the pipeline ended there and was charged its time; if
it found no incumbent (18 and 20 instances), $U=+\infty$. Otherwise OBBT ran with
cutoff $U$, and the final solve started on the tightened box with the incumbent
as a start solution, $U$ as a cutoff or objective limit, and a time limit of
600 seconds minus the root and OBBT time. The control arm was identical but used
the box from feasibility-based tightening. The baselines were the default
solvers on the original model, and for Gurobi also its built-in aggressive OBBT
setting (\texttt{OBBT=3}). A diagnostic arm, kept separate, used the known
optimal value as cutoff, the box of the \texttt{fp} rule, no root run, and no
start solution;
it was run with seed $0$ on the instances that the solver's root run did not
solve. Final solves with identical boxes were shared between rules, each rule
being charged its own OBBT time.

\emph{Metrics and timing.} A run counts as solved if optimality is proved within
600 seconds of total time and the objective is within a relative $10^{-3}$ of the
known value. Times are shifted geometric means with shift one second over
instance--seed pairs, using the total time capped at 600 seconds; node ratios are
shifted geometric means with shift ten over pairs solved by both arms. The
McCormick LP gap closed by a box $B$ is
$\min\{1,\max\{0,(L(B)-L_0)/(f^*-L_0)\}\}$, where $L$ is the bound of the OBBT
relaxation on $B$, $L_0$ its bound on the box from feasibility-based tightening,
and $f^*$ the optimal value; it is averaged over instances with
$f^*-L_0>10^{-6}\max\{1,|f^*|\}$. All runs were wall-clock timed on a machine with
18 cores and 36 hardware threads, with up to 24 concurrent single-thread jobs and
load averages of about 23--25.

\emph{Correction of two solved counts.} Two SCIP runs reported optimality with
the value $-79.35$ instead of $-79.75$: \texttt{crudeoil\_lee1\_05} in the control
arm and \texttt{crudeoil\_lee1\_09} in the \texttt{r5} pipeline, both with seed
$1$. Both are numerical failures of the final solve; the corresponding boxes
contain the reference solution. The solved counts in this appendix exclude
them. All time, node, and final-solve summaries below are the archived
originals: they use the recorded times of these two runs, 30.8 and 105.0
seconds, and count both runs as solved when selecting pairs solved by both
arms. With 678 pairs per arm, the two runs have little weight, but these
summaries are not corrected reanalyses.

\subsection{Solver-level results}

\begin{table}[htbp]
\centering
\small
\begin{tabular}{@{}llrrrrr@{}}
\toprule
Solver & Arm & Solved & Time (s) & Time ratio & Node ratio & Root+OBBT (s)\\
\midrule
Gurobi & default & 640 & 2.37 & 1.00 & 1.00 & 0\\
Gurobi & \texttt{OBBT=3} & 640 & 2.49 & 1.05 & 0.82 & 0\\
Gurobi & control & 642 & 2.57 & 1.08 & 1.05 & 0.5\\
Gurobi & \texttt{r1} & 637 & 3.11 & 1.31 & 0.99 & 14.2\\
Gurobi & \texttt{r5} & 638 & 3.40 & 1.44 & 0.83 & 17.6\\
Gurobi & \texttt{ad0.5} & 637 & 3.14 & 1.33 & 0.97 & 14.5\\
Gurobi & \texttt{ad0.8} & 637 & 3.23 & 1.37 & 0.90 & 15.3\\
Gurobi & \texttt{fp} & 637 & 4.08 & 1.73 & 0.77 & 30.8\\
\addlinespace
SCIP & default & 575 & 7.18 & 1.00 & 1.00 & 0\\
SCIP & control & 573 & 8.30 & 1.16 & 1.11 & 3.8\\
SCIP & \texttt{r1} & 570 & 8.67 & 1.21 & 1.05 & 16.2\\
SCIP & \texttt{r5} & 563 & 8.84 & 1.23 & 0.85 & 19.7\\
SCIP & \texttt{ad0.5} & 570 & 8.71 & 1.21 & 0.98 & 16.5\\
SCIP & \texttt{ad0.8} & 571 & 8.69 & 1.21 & 0.88 & 17.1\\
SCIP & \texttt{fp} & 561 & 9.50 & 1.32 & 0.73 & 31.4\\
\addlinespace
Gurobi & known-value \texttt{fp} & 217/220 & 9.26/4.50 & 2.06 & 0.70 & 48.5\\
SCIP & known-value \texttt{fp} & 220/221 & 15.1/12.3 & 1.23 & 0.55 & 43.1\\
\bottomrule
\end{tabular}
\caption{Precursor study, all 339 instances and two seeds (678 pairs). Ratios
are relative to the default solver; the last column is the mean root and OBBT
time per pair. The known-value diagnostic uses seed $0$ on 238 (Gurobi) and 273
(SCIP) instances and is compared with the default on the same pairs; its time
includes OBBT.}
\label{tab:precursor-solver}
\end{table}

Table~\ref{tab:precursor-solver} gives the solver-level results. No pipeline
solved more pairs than the default, and every pipeline was slower in shifted
geometric mean. Gurobi's built-in OBBT reduced nodes to $0.82$ at a time ratio of
$1.05$ without a restart. Even with the known optimal value as cutoff, the
\texttt{fp} pipeline was slower than the default on the instances it was run on.
On the 69 (Gurobi) and 122 (SCIP) instances that the default needed at least ten
seconds for or did not solve in some seed, the archived summaries give the
pipelines node ratios of
0.99--1.02 and 1.06--1.27 and time ratios of 1.15--1.23 and 1.27--1.51. They
solved 97--98 of these pairs against 100 for the default with Gurobi, and
127--137 against 141 with SCIP.

\begin{table}[htbp]
\centering
\small
\begin{tabular}{@{}llrrrr@{}}
\toprule
Solver & Rule & Time ratio & Node ratio & Final-solve ratio & Solved (control)\\
\midrule
Gurobi & \texttt{r1} & 1.27 & 0.93 & 1.02 & 435 (440)\\
Gurobi & \texttt{ad0.8} & 1.33 & 0.81 & 0.98 & 435 (440)\\
Gurobi & \texttt{fp} & 1.79 & 0.65 & 0.98 & 435 (440)\\
SCIP & \texttt{r1} & 1.05 & 0.96 & 1.02 & 438 (441)\\
SCIP & \texttt{ad0.8} & 1.06 & 0.77 & 1.00 & 439 (441)\\
SCIP & \texttt{fp} & 1.17 & 0.62 & 0.96 & 429 (441)\\
\bottomrule
\end{tabular}
\caption{Effect of the tightened box: pipelines against the control on the
instances where OBBT ran (238 for Gurobi, 273 for SCIP). The final-solve ratio
excludes root and OBBT time and is taken over pairs solved by both arms.}
\label{tab:precursor-control}
\end{table}

Table~\ref{tab:precursor-control} isolates the box. The pipelines and the
control share the root run, the incumbent, and the cutoff, and differ only in the
box. The tightened boxes reduced nodes, but the savings in the final solve did
not offset the root and OBBT time in the total-time summaries. Split by the
McCormick LP gap closed by the \texttt{fp} box, the pairs solved by both arms
show where the node reductions occur. With Gurobi, 112 pairs with closure at
least $0.5$ had node ratio $0.18$ and final-solve ratio $0.87$, at a mean control time of 4.9 seconds; 98 pairs with
positive closure below $0.5$ had $0.77$ and $0.95$ (24 seconds); and 219 pairs
without closure had $0.96$ and $1.00$ (27 seconds). With SCIP the corresponding
values were 142 pairs with $0.20$ and $0.65$ (1.2 seconds), 76 pairs with $0.83$
and $1.02$ (50 seconds), and 198 pairs with $0.97$ and $1.00$ (61 seconds). The
largest mean node reductions occurred in the groups whose mean control time was
only a few seconds.

\subsection{OBBT-level results}

\begin{table}[htbp]
\centering
\small
\begin{tabular}{@{}lrrrrrrr@{}}
\toprule
Cutoff & Instances & 1 & 2 & 3 & 5 & 10 & \texttt{fp}\\
\midrule
Known optimal value & 305 & 0.236 & 0.274 & 0.296 & 0.322 & 0.338 & 0.346\\
Gurobi root incumbent & 235 & 0.147 & 0.181 & 0.203 & 0.231 & 0.247 & 0.250\\
SCIP root incumbent & 267 & 0.176 & 0.213 & 0.237 & 0.260 & 0.278 & 0.287\\
\bottomrule
\end{tabular}
\caption{Mean McCormick LP gap closed after the given number of sequential
rounds and with the rule \texttt{fp}; trajectories that stopped earlier keep
their final value. Shifted geometric means of cumulative OBBT time were 1.1,
1.4, and 1.1 seconds after one round and 2.2, 2.9, and 2.1 seconds with
\texttt{fp}.}
\label{tab:precursor-rounds}
\end{table}

Table~\ref{tab:precursor-rounds} shows diminishing returns from additional
rounds. The median closure was zero for every rule. With the known optimal
value as cutoff, the first round changed at least one bound on 239 of the 339
instances. Of the 104 \texttt{fp} trajectories that ended within their first
round, 97 completed that round without a bound change and seven were
interrupted by the 400-second cap. A further 22 trajectories changed a bound in
the first round and then completed a second round with no accepted change; one
more recorded no change in a second round that the cap interrupted. OBBT time
was heavy-tailed: the median was 0.04 seconds for one round and 0.23 seconds for
\texttt{fp}, but 36 instances needed more than 10 seconds for the first round
and 14 more than 120 seconds, and 18 \texttt{fp} trajectories reached the
400-second cap.

Among the 149 trajectories with at least four rounds in which a variable moved,
the median of $\rho$ over rounds 3--10 exceeded $0.95$ on 79 and was at most $0.5$
on 34. This sample excludes every trajectory that converged quickly, so it does
not show that slow contraction is typical. With a cutoff equal to the optimal
value up to $10^{-6}$, it also cannot separate a stalled limit from slow
approach to a cutoff-dependent floor.

The adaptive rule \texttt{ad0.8} stopped after the first round on $46$--$54\%$
of instances, depending on the cutoff source. Restricting later rounds to moved
variables and their bilinear partners saved only $3\%$ of the LPs (105,878
against 109,107), because these neighborhoods covered almost all nonlinear
variables. Feasibility filtering reduced the LPs of the first round from 87,443
to 25,712 and their time from 7,834 to 5,277 seconds, while the mean closure
changed from $0.239$ to $0.236$. Until the \texttt{fp} rule stopped them,
simultaneous (Jacobi) rounds used 300,710 LPs and 12,053 seconds, against
289,941 LPs and 12,031 seconds for sequential rounds; both had a median of four
rounds and a mean closure of $0.346$.

These filtering and Jacobi--sequential comparison totals are taken from
the archived report; part of their calculation code is absent from the
saved analysis script.

Every one of the 4,573 tightened boxes that belonged to an instance with a
reference solution contained it numerically, with tolerance $10^{-6}(1+|x_j|)$
in each coordinate. Three SCIP runs failed with an LP solver error and count as unsolved.
